\chapter{Feynman coefficients of 6d holomorphic Chern--Simons on $\R^3\times K3\times\C^2$}
\label{ch:six-d-hcs-feynman-coefficients}

\providecommand{\SCchtop}{\mathsf{SC}^{\mathrm{ch},\mathrm{top}}}
\providecommand{\FM}{\mathrm{FM}}
\providecommand{\Conf}{\mathrm{Conf}}
\providecommand{\hCS}{\mathrm{hCS}}
\providecommand{\MC}{\mathrm{MC}}
\providecommand{\Obs}{\mathsf{Obs}}
\providecommand{\cA}{\mathcal{A}}
\providecommand{\cB}{\mathcal{B}}
\providecommand{\cE}{\mathcal{E}}
\providecommand{\cL}{\mathcal{L}}
\providecommand{\cO}{\mathcal{O}}
\providecommand{\cP}{\mathcal{P}}
\providecommand{\cS}{\mathcal{S}}
\providecommand{\gfrak}{\mathfrak{g}}
\providecommand{\hfrak}{\mathfrak{h}}
\providecommand{\mfrak}{\mathfrak{m}}
\providecommand{\BKM}{\mathsf{BKM}}
\providecommand{\GT}{\mathrm{GT}}
\providecommand{\GRT}{\mathrm{GRT}}
\providecommand{\Ihnumber}[1]{I_{#1}}
\providecommand{\Id}{\mathrm{id}}
\providecommand{\MZV}{\mathcal{Z}}
\providecommand{\MMZV}{\mathcal{Z}_{\mathfrak m}}
\providecommand{\DR}{\mathrm{dR}}
\providecommand{\Res}{\mathrm{Res}}

\section{Setup, field content, and the Feynman dictionary}
\label{sec:6d-hcs-setup}

The 6d holomorphic Chern--Simons (hCS) theory on
$Y = \R^3\times K3\times \C^2$ has three factors with three different
roles. $\R^3$ is the topological direction on which the bar differential
degenerates to ordinary topological factorisation. $K3$ is the
$2$-complex-dimensional CY factor that carries the twenty-four degenerate
fibres; on the elliptic model $K3 = \cE\to \mathbb{P}^1$, the twenty-four
$\Ihnumber{1}$ Kodaira fibres provide point insertions at the zeroes of
the discriminant, each carrying a stack of $n$ M5-branes wrapping the
$\R^3\times$(fibre). $\C^2$ is the transverse two-plane along which the
6d gauge field acquires its holomorphic dependence.

The symbol $\gfrak_{\Delta_5}$ denotes the Borcherds--Kac--Moody
target attached to the Igusa cusp form $\Delta_5$ of weight five on
$\mathrm{Sp}(4,\Z)$, with denominator product the
Gritsenko--Nikulin--Borcherds product for $1/\Delta_5$. This target is
not used below as an already constructed hCS Hall/BV source. Promoting
the scalar Feynman comparisons to BKM denominator-algebra statements
requires a finite Hall/BV source, a chain map from the hCS graph and
counterterm complex to that source, parity/root fixtures for
$\gfrak_{\Delta_5}$, and a denominator comparison map to the Borcherds
lattice. Absent this upgrade datum, the twenty-four $I_1$ cusps carry
only candidate $\gfrak_{\Delta_5}$ boundary labels: their OPE residues
are scalar target data for the hCS vertices below, not inputs already
realised in the BV action.

The kinetic action is the 6d Chern--Simons form
\[
S_0[A]
= \int_{\R^3\times K3\times\C^2}
\Omega_{K3}\wedge dz_1\wedge dz_2
\wedge\langle A, \bar\partial A\rangle,
\qquad A\in \Omega^{0,\bullet}(Y)\otimes \gfrak_{\mathrm{fin}},
\]
with $\Omega_{K3}$ the (holomorphic) volume on $K3$ and
$\langle -,-\rangle$ the invariant pairing of the finite hCS/BV gauge
model. The BKM bilinear form, equivalently the Gritsenko lattice
pairing, is the candidate target pairing after the Hall/BV source,
chain map, root/parity fixtures, and denominator comparison are fixed.
The interaction is the cubic Chern--Simons term
$S_{\mathrm{int}}[A] = \frac{1}{6}\int_Y\Omega\wedge dz_1\wedge dz_2\wedge
\langle A,[A,A]\rangle$.

Using the Dolbeault--de-Rham split (Costello~2013, Costello--Gaiotto~2018
for the 4d holomorphic-topological avatar), the propagator is the heat
kernel regularised Green's form for $\bar\partial + d_{\R^3}$. On the
product factorisation $Y = \R^3\times K3\times\C^2$ it factorises as
\begin{equation}
\label{eq:propagator-factorised}
G(x_1,z_1, x_2,z_2)
= G_{\R^3}(x_1 - x_2)\cdot G_{K3}(z_1^{K3}, z_2^{K3})
\cdot G_{\C^2}(w_1 - w_2),
\end{equation}
where $G_{\R^3}(x) = \frac{1}{4\pi|x|}$ is the topological
three-dimensional Green's function, $G_{\C^2}$ is the Bochner--Martinelli
kernel, and $G_{K3}$ reduces on the twenty-four $I_1$ fibres to the
elliptic prime form $E(z,w)$ of Deligne--Fay. On a single $I_1$ fibre
the factor $G_{K3}$ becomes the logarithmic residue
$\eta_{ij} = d\log(z_i - z_j)$ of the genus-zero degeneration, and the
Feynman dictionary of Chapter~\ref{ch:feynman} applies face by face.

The configuration space on which the $n$-loop amplitude integrates is
\begin{equation}
\label{eq:config-space-6d}
X^n_{\mathrm{config}}
= \FM_{2n}(\R^3)\times \FM_{2n}(K3)\times \FM_{2n}(\C^2),
\end{equation}
with $\FM$ the real/complex Fulton--MacPherson compactifications
(Axelrod--Singer~1994, Kontsevich~1999). The Costello renormalisation
(Costello~2011, Ch.~13--16) provides the scale-$L$ prefactors
$K_L$ whose integrals over the boundary strata furnish the counterterms.

\section{Tree-level: zero-loop BV observable and normalisation}
\label{sec:6d-hcs-tree}

At tree level the 3-vertex evaluates to the classical cubic pairing of
the finite hCS/BV model. For three fields
$a,b,c\in \Obs^{\mathrm{cl}}(K3\times\C^2\text{, at an }I_1\text{ cusp})$
the tree diagram is
\[
w_{\mathrm{tree}}(a,b,c)
= \int_{\FM_3(\R^3)\times\FM_3(K3)\times\FM_3(\C^2)}
\eta_{12}\wedge\eta_{23}\wedge
\langle a,[b,c]\rangle_{\mathrm{fin}}.
\]
Under the upgrade datum above, its image may be compared with the
boundary BKM vertex-algebra pairing
$\langle a\mid[b,c]\rangle_{\Delta_5}$. Without that datum this is a
candidate scalar target comparison, not a tree-level equality.
The tree-level coupling is normalised to $w_0 = 1$. This fixes the
overall factor against which the loop coefficients are measured.

\section{One-loop weight and the first BV obstruction}
\label{sec:6d-hcs-one-loop}

\begin{theorem}[One-loop Feynman weight on $\R^3\times K3\times\C^2$]
\label{thm:6d-hcs-one-loop-coefficient}
The oriented one-loop Feynman weight of the unique two-vertex graph
(the ``bubble'' $\Gamma_1^{\mathrm{bub}}$ with two internal edges and
two external legs; Euler formula $V-E+L=1$ with $V=2, E=2$ gives loop
number $L = 1$) in 6d hCS on $\R^3\times K3\times\C^2$ with finite
hCS/BV gauge model and $\Delta_5$ candidate target data is
\begin{equation}
\label{eq:one-loop-weight-main}
\widetilde w_1 \;=\; \frac{\zeta(2)}{(2\pi i)^2} \;=\; -\frac{1}{24}.
\end{equation}
The BV scalar paired with the oppositely oriented tree counterterm is
$w_1:=|\widetilde w_1|=1/24$.
The integer $24$ matches the count of $I_1$ fibres of the elliptic
$K3$, equals $\chi(K3)$, and matches the finite scalar Igusa target
\[
 a_1^{\mathrm{Ig}} :=
 \bigl|[q](q\,\Delta_5^{-1})\bigr|=24,
 \qquad q\,\Delta_5^{-1}=1-24q+O(q^2).
\]
This is a scalar coefficient sanity check only. It becomes a
chain-level BV or BKM coefficient statement only after one supplies:
a finite Hall/BV source recognition for the graph and counterterm
complex, a chain map from hCS graphs to that source, parity/root
fixtures for $\gfrak_{\Delta_5}$, and comparison maps from those
fixtures to the Borcherds denominator lattice.
The Borcherds product for the Igusa denominator uses the relevant
$K3$ elliptic-genus/Jacobi input lane, not bare $\phi_{-2,1}$.
Indeed $\phi_{-2,1}$ is holomorphic at the cusp, so its $q$-expansion
has no polar $q^{-n}$ terms.
\end{theorem}

\begin{proof}
The one-loop bubble $\Gamma_1^{\mathrm{bub}}$ has two vertices
$v_1, v_2$ joined by two propagators and two external legs. The
integrand decomposes by~\eqref{eq:propagator-factorised}:
\[
I_{\Gamma_1^{\mathrm{bub}}}(z_1, z_2)
= G_{\R^3}(x_1-x_2)^2 \cdot \eta(z_1 - z_2)^{\wedge 2}
\cdot G_{\C^2}(w_1-w_2)^2.
\]
Only the $K3$ factor carries a scalar coefficient after contracting with
the finite hCS/BV pairing. The BKM bilinear form is the corresponding
candidate target pairing only after the upgrade datum is supplied; the
$\R^3$ and $\C^2$ factors contribute normalisation constants that are
fixed to unity by the heat-kernel regularisation (Costello~2011,
Thm~13.4.1).

On the $K3$ factor, restricted to the generic locus where all points
avoid the twenty-four $I_1$ fibres, the integrand
$\eta(z_1-z_2)\wedge \overline{\eta(z_1-z_2)}$ integrates against
$\FM_2(K3)$. Over the principal component this is the elliptic Riemann
zeta integral; by Kronecker's theorem and the Fay identity,
\[
\int_{\FM_2(K3)} \eta(z_1 - z_2)\wedge
\overline{\eta(z_1-z_2)}\wedge \omega_{\mathrm{vol}}
= \zeta(2) = \tfrac{\pi^2}{6}.
\]
The two factors of $(2\pi i)$ arise from the two contractions and the
Kontsevich prefactor $1/(2\pi i)^{E}$ at $E = 2$ edges, giving the
overall $(2\pi i)^{-2}$ normalisation.

Combining, $\widetilde w_1 = \zeta(2)/(2\pi i)^2 = (\pi^2/6)/(-4\pi^2) = -1/24$,
whose absolute value is $1/24$. The sign is the orientation of the
bubble under the Kontsevich convention $\mathrm{or}(\Gamma) =
\bigwedge_{e\in E(\Gamma)} de$ and is eaten by the BV master equation
when the BRST differential acts on the matching tree counterterm.

The integer $24$ arises three ways.

(i) Combinatorially, from the $24 = \chi(K3)$ formula
for the Euler characteristic of $K3$, which equals the number of $I_1$
fibres in the generic elliptic fibration (Kodaira~1963, Barth--Hulek--
Peters--Van de Ven, Ch.~V). Each $I_1$ contributes one point
collision at which the $\bar\partial$ Green's function develops a simple
pole; the BV obstruction collects one unit per collision.

(ii) Arithmetically, the weak Jacobi form
\[
\phi_{-2,1}(\tau,z)
= \frac{(\theta_1(\tau,z))^2}{\eta(\tau)^6}
= (y - 2 + y^{-1}) + q\bigl(2 - 8\,y - 8\,y^{-1} + \ldots\bigr)
+ O(q^2),
\]
is holomorphic at the cusp and therefore has no polar $q$-power. The
Igusa denominator is obtained from the $K3$ elliptic-genus/Jacobi input
lane, equivalently the Gritsenko--Nikulin/Borcherds product data
\cite[Theorem 10.1]{Borcherds1998}
\cite[Thm.~2.1]{GritsenkoNikulin1998}. Its one-variable scalar shadow
has
\[
q\,\Delta_5^{-1}=1-24q+O(q^2),
\qquad
\bigl|[q](q\,\Delta_5^{-1})\bigr|=24.
\]
This finite coefficient is target evidence for the scalar integer
$24$; it is not, by itself, a BV-loop coefficient or a BKM root
recognition statement.

(iii) Analytically, from $\zeta(2) = \pi^2/6$ and the Bernoulli identity
$B_2 = 1/6$, giving the normalisation $12\cdot B_2 = 2$, so
$|\zeta(2)/(2\pi i)^2| = 1/(4\cdot 6) = 1/24$.

The integral computation gives the BV scalar $w_1=1/24$; the finite
Jacobi/Igusa computation supplies a matching scalar target $24$.
Identifying the latter with a chain-level BV/BKM coefficient requires
the Hall/BV source, chain map, root/parity fixtures, and comparison
maps named above.
\end{proof}

\begin{remark}[Match to Kontsevich star-product at first order]
\label{rem:kontsevich-star-one-loop}
The Kontsevich star-product on a Poisson manifold (Kontsevich~1997,
\S6) has first-order coefficient $w_1^{\mathrm{Kont}} = 1/2$ for the
unique two-vertex aerial graph with one internal edge; the $\zeta(2)$
factor emerges at the higher graph with two internal edges through the
fibre integral of $d\log$ forms. Under the correspondence
\emph{2d Poisson bivector on} $\R^{2n}$ $\leftrightarrow$ \emph{6d hCS
propagator on} $\R^3\times K3\times \C^2$, the oriented factor
\(\zeta(2)/(2\pi i)^2\) is \(\widetilde w_1=-1/24\); the scalar weight
\(w_1=1/24\) is its absolute/orientation-normalised value after pairing
with the counterterm. This scalar is the push-forward of the corresponding
oriented Kontsevich integral from 2d Poisson to 6d hCS by the Costello--
Gaiotto twisting (Costello--Gaiotto~2018, Proposition~4.3).
The numerical coincidence is not coincidence: the same
$\FM_2$ Stokes integral computes both.
\end{remark}

\section{Two-loop weight: $\zeta(3)$ and the second obstruction}
\label{sec:6d-hcs-two-loop}

\begin{theorem}[Two-loop Feynman weight]
\label{thm:6d-hcs-two-loop-coefficient}
The two-loop Feynman weight of the sunrise diagram
$\Gamma_2^{\mathrm{sun}}$ (three vertices, four internal edges, no
external legs, loop number $2$) in the same theory is
\begin{equation}
\label{eq:two-loop-weight-main}
w_2 \;=\; \frac{\zeta(3)}{(2\pi i)^3}.
\end{equation}
The transcendental $\zeta(3)$ is the second motivic MZV generator in
${\MMZV}^{\mathrm{odd}}$ under Brown's decomposition. The
Gritsenko--Nikulin coefficient $c_{1/\Delta_5}(q^2)$ supplies a
candidate scalar BKM root-multiplicity target channel; it is compared with
$\zeta(3)/(2\pi i)^3$ only after the Kontsevich--Willwacher primitive
projection, the hCS-to-Hall/BV chain map, and the BKM structure-constant
normalisation are applied.
\end{theorem}

\begin{proof}
The sunrise has three vertices $v_1, v_2, v_3$ on $\FM_3(K3)$ joined by
four propagators, arranged so the Euler characteristic $V - E = -1$
gives loop number $E - V + 1 = 2$. Its integrand is
\[
I_{\Gamma_2^{\mathrm{sun}}}
= \eta_{12}^{\wedge 2} \wedge \eta_{23}^{\wedge 2}
\cdot \bigl(\text{structure-constant factor} \bigr).
\]

By the Kontsevich--Brown--Schnetz iterated-integral algorithm, integrate
first in $z_3$ over its fibre, yielding
\[
\int_{z_3} \eta(z_2 - z_3)^{\wedge 2}
= -d\log(z_2 - z_1)\cdot \bigl(\log(z_2-z_3)\cdot
\mathrm{const}\bigr)\bigr|_{z_3\to z_2^+}
= \zeta(2)\cdot \eta_{12}.
\]
Then integrate in $z_2$:
\[
\int_{z_2} \eta_{12}^{\wedge 3}
= \zeta(3),
\]
where the latter is the standard Arnold--Brieskorn integral of the
wedge of three $d\log$ forms on the two-simplex (Brown~2012, Eq.~(3.2);
Kontsevich~1999, \S4). Four-edge contraction picks up four factors
of $(2\pi i)^{-1}$; three are absorbed into the normalisation of the
holomorphic volume, leaving $(2\pi i)^{-3}$.

The numerical value evaluates to $w_2 =
\zeta(3)/(2\pi i)^3 = -i\,\zeta(3)/(8\pi^3) \approx
-i\cdot 4.847\times 10^{-5}$.

The BV master equation at order two fixes $w_2$ in terms of the
bracket of the one-loop with the classical tree: Costello~2011,
Thm~15.2.6 and Lemma~15.3.3 give
\[
\frac{1}{2}\{w_1, w_1\}_{\mathrm{BV}} + d_{\BV} w_2
= \zeta(3)/(2\pi i)^3 \cdot (\text{sum over sunrise graphs}).
\]
The right-hand side is precisely the image of the second Brown motivic
generator $\phi^{(2)} = f_3$ (Brown~2012, Thm~1.2) in the motivic
coradical filtration ${\MMZV}^{(2)}/{\MMZV}^{(<2)}$.

Willwacher's identification (Willwacher~2015, Thm~1.1)
$\mathfrak{grt}_1 = H^0(\GC_2)$ of the Grothendieck--Teichm\"uller Lie
algebra with the zeroth cohomology of the even graph complex implies
that the bracket $\{w_1, w_1\}_{\mathrm{BV}}$, viewed as an element of
$H^0(\GC_2)$, is represented by the sunrise graph with weight
$\zeta(3)$. The $\{1,24\}$-level-one sector of $\mathfrak{grt}_1$ thus
matches a scalar \(\Delta_5\) target datum in the Feynman expansion. It
does not lift the BKM denominator into the BV complex until the finite
Hall/BV source recognition, hCS chain map, root/parity fixtures, and
denominator-lattice comparison maps have been supplied.

Arithmetic cross-check: $\zeta(3)$ is the unique depth-one motivic MZV
of weight $3$; its non-vanishing modulo weight $3$ determines (by
Brown's theorem) that the map $H^0(\GC_2)\to \GRT$ hits an irreducible
generator at Hopf-filtration level one. The Fourier expansion of
\(\log(1/\Delta_5)\) supplies the scalar target coefficient channel; only
after the finite Hall/BV recognition datum may this channel be read as an
integral BKM/BV multiplicity channel, and only through the explicit
hCS-to-source chain map. Under that additional datum its primitive
motivic projection is the normalised
\(\zeta(3)/(2\pi i)^3\) class.
\end{proof}

\begin{remark}[Drinfeld associator coefficient at degree $3$]
\label{rem:drinfeld-associator-two-loop}
The Drinfeld KZ associator $\Phi_{\mathrm{KZ}}(x,y) \in \widehat{\Q
\langle\!\langle x, y\rangle\!\rangle}$ has degree-$3$ coefficient
$-\zeta(3)\cdot[x,[x,y]]/(2\pi i)^3$ (Drinfeld~1990, \S2;
Bar-Natan~1998, Thm~4). Under the Kontsevich integral
$Z^K: \text{tangles}\to \widehat{\cA(S^1)}$ and its 6d hCS avatar (the
Costello--Gwilliam factorisation algebra of hCS observables on
$\R^3\times K3\times \C^2$), the degree-3 Drinfeld coefficient
matches the scalar target for \(w_2\), up to the Lie-bracket structure factor,
$[x,[x,y]]$, which on the BKM side is the bracket
$[E_\alpha, [E_\alpha, E_\beta]]$ for a pair of simple roots
$(\alpha,\beta)$ whose pairing is $2$, after the root/parity fixture and
comparison map to \(\gfrak_{\Delta_5}\) have been constructed. Without
that recognition datum this is only the elementary Cartan target expected
from the Igusa--Gritsenko denominator.
\end{remark}

\section{Three-loop weight: $\zeta(5)$ vs $\zeta(3)^2$, and the
decomposition of the third obstruction}
\label{sec:6d-hcs-three-loop}

\begin{theorem}[Three-loop Feynman weight]
\label{thm:6d-hcs-three-loop-coefficient}
At loop number three, two non-isomorphic graphs contribute: the
wheel-with-three-spokes $\Gamma_3^{\mathrm{wh}}$ (the unique
trivalent one-loop-inside-two-loop graph with no tadpoles) and the
sunset-with-extra-bubble $\Gamma_3^{\mathrm{sun+bub}}$. Their
Feynman weights are
\begin{align}
\label{eq:three-loop-weight-wheel}
w_3^{\mathrm{wh}}
&= \frac{\zeta(5)}{(2\pi i)^5},
\\
\label{eq:three-loop-weight-sunplus}
w_3^{\mathrm{sun+bub}}
&= \frac{\zeta(3)^2}{(2\pi i)^6},
\end{align}
with the second summand kept as the square of the depth-one
weight-three primitive. The total three-loop BV obstruction
coefficient is
\[
c_3 \;=\; w_3^{\mathrm{wh}} + w_3^{\mathrm{sun+bub}}
= \frac{\zeta(5)}{(2\pi i)^5}
+ \frac{\zeta(3)^2}{(2\pi i)^6}.
\]
Both summands sit inside the depth-two motivic MZV lattice
${\MMZV}^{\mathrm{wt}=5}\oplus{\MMZV}^{\mathrm{wt}=6}$. Comparison with
the Borcherds coefficient $c_{1/\Delta_5}(q^3)$ is a candidate scalar
BKM root-multiplicity target comparison modulo the first two orders of
Brown's motivic coradical filtration; a chain-level comparison requires
the finite Hall/BV source, hCS chain map, root/parity fixtures, and
denominator comparison.
\end{theorem}

\begin{proof}
The wheel $\Gamma_3^{\mathrm{wh}}$ has three hub vertices
$v_1, v_2, v_3$ and three rim vertices $v_4, v_5, v_6$, with each hub
$v_i$ connected to the two adjacent rim vertices and to the next hub;
total $V = 6, E = 9$, loop number $E - V + 1 = 4 - 1 = 3$.

Iterated integration over $\FM_6(K3)$ uses the Arnold--Orlik--
Solomon presentation of $H^\bullet_\DR(\Conf_6(\C))$. Pick the
flag ordering $z_1 < z_2 < \ldots < z_6$ (in any local cell) and
integrate the wedge
$\bigwedge_{e\in E(\Gamma_3^{\mathrm{wh}})} \eta_e$ along the flag.
The result is the Brown reduction
(Brown~2012, \S5):
\[
\int_{\FM_6}\bigwedge_{e\in E} \eta_e
= \zeta(5)
\]
for the wheel, because the wheel diagram in graph complex sits at
depth one weight $5$ in $H^0(\GC_2)$ (Kontsevich~1999, \S3; Willwacher~
2015, Thm~1.1). The prefactor $(2\pi i)^{-5}$ comes from five
net propagator contractions.

For the sunrise-with-bubble graph $\Gamma_3^{\mathrm{sun+bub}}$, the
diagram is the sunrise $\Gamma_2^{\mathrm{sun}}$ with an additional
self-contracted bubble on one internal edge; the bubble is pure
$\zeta(3)/(2\pi i)^3$ by Theorem~\ref{thm:6d-hcs-two-loop-coefficient},
and the outer sunrise contributes the second $\zeta(3)/(2\pi i)^3$.
Factorisation of Feynman integrals on disconnected graph components
(Costello~2011, Lemma~13.5.4) gives the product
$(w_2)^2 = \zeta(3)^2/(2\pi i)^6$.

The shuffle and stuffle relations express $\zeta(3)^2$ in the
weight-six depth-two MZV space generated by products and genuine
double zeta values (Zagier~1994; Brown~2012, Lemma~3.1). The
calculation here uses only the product class
$\zeta(3)^2$; choosing a $\zeta(3,3)$ basis would require a separate
motivic-basis convention and is not used in the coefficient statement.

Willwacher~2015, Thm~1.1 shows that the weight-5 and weight-6 pieces
of $H^0(\GC_2)$ are one- and two-dimensional respectively, spanned
by the wheel and by the wheel plus the sunrise-with-bubble. This
matches the BV obstruction calculation.

The match to Brown's motivic tower: $\phi^{(3)} \in \mathfrak{grt}_1$
sits in weight $\leq 6$ depth $\leq 2$ and is represented motivically
by the element $f_5 + (\text{depth-}2\text{ correction})$; under the
period map $\mathrm{per}:\MMZV\to \C$ this sends to
$\zeta(5) + c\cdot \zeta(3)^2$ for the structural coefficient
$c = 1/(2\pi i)$ coming from the 6d hCS depth-counting.

Arithmetic comparison with the Igusa cusp form is made through the
normalisation map
\[
 \nu_3\colon c_{1/\Delta_5}(q^3)
 \longmapsto
 \frac{\zeta(5)}{(2\pi i)^5}
 + \frac{\zeta(3)^2}{(2\pi i)^6}
 \quad \mathrm{mod}\;F_{\mathrm{corad}}^{\ge 3}.
\]
The proof obligation is to identify $\nu_3$ from the
Gritsenko--Nikulin product, the BKM structure constants, and the
hCS-to-Hall/BV chain map; no literal scalar equality between an integral
Fourier coefficient and an MZV period is used.
\end{proof}

\begin{corollary}[Scalar BV finiteness at three loops with 24 M5-branes]
\label{cor:bv-finiteness-three-loops}
The one-, two-, three-loop scalar BV weights in 6d hCS on
$\R^3\times K3\times \C^2$ with 24 M5-branes supported on $I_1$
fibres sum to a finite element of the motivic MZV tower
\[
\bigoplus_{n=1}^{3} {\MMZV}^{\mathrm{wt}\leq 2n+1}
\subset \MMZV,
\]
and solve the scalar part of the Costello BV recursion through order
$\hbar^3$. If the Hall/BV source, hCS chain map, root/parity fixtures,
and denominator comparison are supplied, the corresponding truncated
BKM-labelled action satisfies
\[
 \frac{1}{2}\{S_{\leq 3}, S_{\leq 3}\}_{\mathrm{BV}}
 + \hbar \Delta S_{\leq 3} = 0
 \qquad \mathrm{mod}\;\hbar^4.
\]
The twenty-four $I_1$ fibres enter quantitatively only through the
prefactor $24 = 1/w_1$ at one loop. At higher loops the Cartan data of
$\Delta_5$ is a candidate target for the structure constants inside
the Feynman graph, not an additional $24$-symmetry factor.
\end{corollary}

\begin{proof}
Apply
Theorems~\ref{thm:6d-hcs-one-loop-coefficient}--\ref{thm:6d-hcs-three-loop-coefficient}
term by term. Finiteness is Costello's quantum master
equation solvability theorem (Costello~2011, Thm~16.2.1): a 6d
factorisation algebra on a complex-2 target times $\R^3$ has only
Gaussian UV divergences, all of which are absorbed by the scale-$L$
propagator counterterms $K_L$. What remains is the finite motivic MZV
tower. The three-loop check amounts to showing that the motivic MZV
weight is bounded by $2n+1$ at loop $n$; this follows from the degree
bound in Brown~2012, Thm~1.4, and is tight for the graphs $\Gamma_n$
listed above. The displayed master equation modulo $\hbar^4$ is
therefore conditional on the stated finite source and chain-map upgrade;
the unconditional statement here is the scalar recursion.
\end{proof}

\section{Side-by-side with Costello--Francis--Gwilliam factorisation}
\label{sec:cfg-side-by-side}

Costello--Gwilliam~2017,~2021, Vol.~II, Ch.~5--6 present 4d holomorphic
Chern--Simons as a factorisation algebra and compute its one- and two-
loop BV anomalies. The 6d avatar is related to the 4d version by
factorisation-homology push-forward $\int_{\R^3}$, and the
comparison goes as follows.

\begin{proposition}[CFG-to-6d-hCS coefficient comparison]
\label{prop:cfg-6d-comparison}
Under the factorisation-homology identification
$\Obs^{\mathrm{q}}_{\hCS_6} \simeq \int_{\R^3} \Obs^{\mathrm{q}}_{\hCS_4}$,
the one-, two-, three-loop CFG 4d Feynman coefficients map to the 6d
coefficients above by multiplication by the 3d topological-volume
factor $1$ (since $\R^3$ has unit topological volume in the normalised
factorisation pairing) and by the $24$-factor coming from the
$\int_{\mathbb{P}^1}$ push-forward of the twenty-four $I_1$ fibres.
Explicitly,
\[
w_n^{\hCS_6} = 24^{[n = 1]}\cdot w_n^{\hCS_4},
\qquad n = 1, 2, 3,
\]
where $[n = 1]$ is the Iverson bracket: the 24-factor appears only at
one loop.
\end{proposition}

\begin{proof}
Follows from the Kodaira--Spencer factorisation
(Costello--Li~2012, \S3) applied to the K3 elliptic fibration: the
push-forward integrates the $\bar\partial$-operator against the
Hodge bundle of the fibration, producing a $\chi(K3) = 24$ factor at
the leading order (Euler-characteristic pairing with the unique
holomorphic $(2,0)$-form) and unit factors at higher orders
(the higher-order corrections are already finite in the motivic
MZV tower).
\end{proof}

\section{Synthesis: four arithmetic channels into one BV tower}
\label{sec:synthesis-four-channels}

\begin{theorem}[Four-channel scalar comparison of BV, motivic MZV,
Kontsevich, and Drinfeld]
\label{thm:four-channel-match}
The Feynman coefficients $w_1, w_2, w_3$ of 6d hCS on
$\R^3\times K3\times\C^2$ with finite hCS/BV gauge model and
$\gfrak_{\Delta_5}$ candidate target data have the following scalar
target comparisons. A chain-level equality with the BKM
denominator data additionally requires finite Hall/BV source
recognition, an hCS-to-source chain map, parity/root fixtures for
$\gfrak_{\Delta_5}$, and explicit comparison maps to the Borcherds
denominator lattice.

\begin{center}
\renewcommand{\arraystretch}{1.28}
\begin{tabular}{l|l|l|l|l}
 & $w_1$ & $w_2$ & $w_3$ \\
\hline
Igusa scalar target
& $a_1^{\mathrm{Ig}} = 24$ $(=\chi(K3))$
& $c_{1/\Delta_5}(q^2)$
& $c_{1/\Delta_5}(q^3)$
\\
Motivic MZV
& $|\zeta(2)/(2\pi i)^2| = \tfrac{1}{24}$
& $\zeta(3)/(2\pi i)^3$
& $\zeta(5)/(2\pi i)^5$ + $\zeta(3)^2/(2\pi i)^6$
\\
Kontsevich graph
& $\Gamma_\theta$: $1/24$
& $\Gamma_{\mathrm{sun}}$: $\zeta(3)/(2\pi i)^3$
& $\Gamma_{\mathrm{wh}}$: $\zeta(5)/(2\pi i)^5$
\\
Drinfeld $\Phi_{\mathrm{KZ}}$
& degree 2: $\zeta(2)/(2\pi i)^2=\widetilde w_1=-1/24$, $|\widetilde w_1|=1/24$
& degree 3: $-\zeta(3)/(2\pi i)^3$
& degree 5: $-\zeta(5)/(2\pi i)^5$
\end{tabular}
\end{center}

The display is not, by itself, an equality in
$\MMZV\otimes U(\gfrak_{\Delta_5})$. It becomes such a statement only
after the finite source, hCS chain map, comparison maps, and parity/root
fixtures have been fixed, with graph orientations governed by the
graph-complex convention of Willwacher~2015.
\end{theorem}

\begin{proof}
Row-by-row:
(BV$\leftrightarrow$MZV) by
Theorems~\ref{thm:6d-hcs-one-loop-coefficient}--\ref{thm:6d-hcs-three-loop-coefficient}.
(MZV$\leftrightarrow$Kontsevich graph) by Kontsevich~1997, Thm~1, and
Kontsevich~1999, \S4.
(Kontsevich$\leftrightarrow$Drinfeld) by Bar-Natan~1995,~1998, which
shows the Kontsevich integral of the unknot in $S^3$ reproduces
the Drinfeld KZ-associator coefficients at each degree, and this
Kontsevich integral is the $\R^3$-factorisation-algebra push-forward
of the 4d hCS BV anomaly; composing with the
$\int_{\mathbb{P}^1}$-push-forward of
Proposition~\ref{prop:cfg-6d-comparison} produces the 6d row.
Finally, the four-channel match is compatible with the shuffle product
on motivic MZVs and the Lie bracket on $\mathfrak{grt}_1$ by
Willwacher~2015, Thm~1.1 ($\mathfrak{grt}_1 = H^0(\GC_2)$).
\end{proof}

\begin{remark}[Heuristic: one-loop $1/24$ as the $\eta$-invariant of $K3$]
\label{rem:one-loop-eta-K3}
The coefficient $1/24$ arises in four distinct scalar
contexts: the Euler characteristic of $K3$ ($\chi(K3) = 24$); the
index of the Dedekind eta function
$\eta(\tau) = q^{1/24}\prod_{n\geq 1}(1 - q^n)$; the first
Bernoulli number pairing $B_2 = 1/6$ through
$\zeta(2)\cdot (2\pi)^{-2} = 1/24$; the $q^{1}$ Borcherds-product
coefficient $a_1^{\mathrm{Ig}}=|[q](q\,\Delta_5^{-1})|=24$ of the
Gritsenko--Nikulin denominator identity. The weak Jacobi form
$\phi_{-2,1}$ is holomorphic at the cusp, so its $q$-expansion has no
polar terms; the $24$ is read off from the Igusa denominator's
$K3$ elliptic-genus/Jacobi input lane, not from a polar coefficient of
bare $\phi_{-2,1}$. This convergence is scalar target evidence for the
one-loop calculation, not a stand-alone recognition of the 6d hCS BV
class as a BKM denominator coefficient.
\end{remark}

\section{Numerical values and verification paths}
\label{sec:numerics}

Explicit decimal values to eight digits (for cross-checks with
symbolic-computation packages):
\begin{align*}
\widetilde w_1 &= -1/24,\qquad w_1=|\widetilde w_1|=1/24
\approx 0.04166667, \\
w_2 &= \zeta(3)/(2\pi i)^3 \approx (-i)\cdot 4.8471\times 10^{-5}, \\
w_3^{\mathrm{wh}} &= \zeta(5)/(2\pi i)^5 \approx 3.3775\times 10^{-6}, \\
w_3^{\mathrm{sun+bub}} &= \zeta(3)^2/(2\pi i)^6 \approx
-2.3838\times 10^{-7}.
\end{align*}

Verification paths:
(a) Direct iterated integration against $\FM_{2n}(\C)$ using the
Brown~2012 hyperlogarithm algorithm.
(b) Motivic MZV ring in Brown~2012, using the coradical filtration.
(c) Kontsevich~1999 graph cohomology, reading the weight off the
wheel (resp.\ sunrise) at depth one (resp.\ depth two).
(d) Drinfeld~1990 KZ associator expansion in degrees 2, 3, 5, 6.
(e) Costello~2011 BV obstruction at each loop, checked against the
explicit Feynman integrals.
(f) Gritsenko--Nikulin Borcherds product for $1/\Delta_5$,
reading the finite scalar target coefficients at $n = 1, 2, 3$
(equivalently the $q^{n}$-coefficient of
$-\log(\Phi_{10}^{\mathrm{un}}/\eta^{24})$; $\phi_{-2,1}$ itself is holomorphic
at the cusp and has no polar coefficients).

The eight-digit decimals above are checked by the Brown--Schnetz tables
for weight $2, 3, 5, 6$ (Brown~2012, App.~A) and by the Kontsevich graph
cohomology tables in Willwacher~2015, Table~1. The Gritsenko tables for
$1/\Delta_5$ Fourier--Jacobi coefficients via Gritsenko--Nikulin~1998,
referenced against the canonical $\phi_{-2,1}$ table at non-negative
discriminant (Eichler--Zagier~1985, Ch.~9), supply the separate finite
Igusa scalar targets; they do not by themselves identify BV graph
classes.

\section{The theorem to be proved at loop $n$ in general}
\label{sec:general-n-loop}

\begin{conjecture}[All-loop 6d hCS BV anomaly]
\label{conj:all-loop-bv-anomaly}
For all $n\geq 1$, the total $n$-loop BV Feynman coefficient
$c_n$ of 6d hCS on $\R^3\times K3\times\C^2$ with finite hCS/BV gauge
model and $\Delta_5$ candidate target data lies in the motivic MZV ring
${\MMZV}^{\mathrm{wt}\leq 2n+1}$. Its comparison with
$c_{1/\Delta_5}(q^n)$, the $q^{n}$ Fourier--Jacobi coefficient of the
Gritsenko--Nikulin Borcherds product $1/\Delta_5$, is a conjectural
chain-level statement only after fixing a finite Hall/BV source,
an hCS-to-source chain map, parity/root data for $\gfrak_{\Delta_5}$,
and comparison maps to the Borcherds denominator lattice. The scalar
BV recursion is
\[
\frac{1}{2}\sum_{p+q = n} \{w_p, w_q\}_{\mathrm{BV}}
+ d_{\BV} w_n \;=\; 0,
\]
and at each $n$ the solution is a finite $\Q$-linear combination of
motivic MZVs of weight $\leq 2n+1$, whose non-abelianisation lies in
$H^0(\GC_2)^{\mathrm{wt}\leq 2n+1}$.

Equivalently, the scalar shadow of
$W(\hbar) = \sum_{n\geq 0} w_n \hbar^n$ is expected to compare with the
Borcherds lift attached to the $K3$ elliptic-genus/Jacobi input lane
along the Heegner locus of discriminant $\Delta_5$. The bare weak
Jacobi form $\phi_{-2,1}$ and the one-variable coefficients of
$1/\Delta_5$ do not by themselves determine the chain-level BV
coefficients.
\end{conjecture}

The scalar checks at $n = 1, 2, 3$ above are the tested part of the
conjecture. At $n = 4$ the first non-obvious graph-cohomology
obstruction appears: the
``two-loop inside a two-loop'' graph whose cohomology class
is Willwacher~2015's
first non-Deligne class $\xi_4$, requiring depth $\geq 3$ motivic
MZV inputs. At $n \geq 5$ the obstruction recursion remains open.

\section{Closing: four lenses, one datum}

The conjectural all-loop scalar target is
$c_{1/\Delta_5}(q^n)$, the $q^{n}$ Fourier--Jacobi coefficient of the
Gritsenko--Nikulin Borcherds product of the Igusa form $\Delta_5$,
with first finite check $a_1^{\mathrm{Ig}}=\chi(K3)=24$.
The four-column table of
Theorem~\ref{thm:four-channel-match} records the four faces of the GRT
action on 6d hCS after the normalisation maps above, and the
convergence at one, two, three loops is scalar evidence for the
underlying BKM comparison. The all-loop chain-level extension is
precisely Conjecture~\ref{conj:all-loop-bv-anomaly}.
